\section{Delay Without Lying}
\label{sec:clocks-delay}

There are three ways to wait on a Cortex-M and this project has been wrong about all three.

\subsection{asm::delay Is Not a Cycle Count}

\begin{minted}{rust}
// "1000 cycles", according to the name and to every example on the internet.
cortex_m::asm::delay(1000);
\end{minted}

Measured against the cycle counter, one iteration of the loop that \codeinline{asm::delay} expands
to takes about \textbf{three} cycles on this part with this compiler. So
\codeinline{delay(1\_000\_000)} waits three million cycles, not one — a 3× error, and the reason a
2.7 kHz tone computed with it came out as something else entirely.

Two consequences, both practical:

\begin{itemize}
    \item \textbf{Never derive a period from it.} A tone, a baud rate or a bus half-period computed
          with \codeinline{asm::delay} is wrong by a factor nobody wrote down.
    \item \textbf{Do not multiply by three and call it fixed.} The factor is a property of the
          compiler's code generation, not of the hardware; it changed between optimisation levels in
          this project.
\end{itemize}

\subsection{SysTick: Real Milliseconds}

SysTick counts CPU cycles and can therefore produce real time:

\begin{minted}{rust}
pub fn delay_us(&mut self, us: u32) {
    let cycles = (u64::from(us) * u64::from(self.sysclk_hz)) / 1_000_000;
    // SysTick's reload is 24 bits: split anything larger into chunks.
    let mut remaining = cycles;
    while remaining > 0 {
        let chunk = remaining.min(0x00FF_FFFF);
        self.syst.set_reload(chunk as u32);
        self.syst.clear_current();
        self.syst.enable_counter();
        while !self.syst.has_wrapped() {}
        self.syst.disable_counter();
        remaining -= chunk;
    }
}
\end{minted}

It has one owner — the board — and it must never be shared with a scheduler that also wants it. In
this project that rule is enforced by the type system in the smallest possible way: \codeinline{SYST}
is taken once by \codeinline{Board::take()}, and \codeinline{board.delay()} hands it out.

\subsection{Ever the Millisecond: Do Not Divide Early}

A defect worth its own paragraph, because it looks like arithmetic hygiene:

\begin{minted}{rust}
// The sample interval arrives as microseconds: 2500 us, i.e. 416 Hz.
let dt_ms = dt_us / 1000;        // 2 ms
// Every timeout in the state machine is now 20% short.
\end{minted}

Integer division by 1000 truncated 2500 to 2, and every threshold derived from it inherited a 20%
error. The classifier's free-fall window, its landing window and its settle timer were all wrong at
once, and every one of them was measured against real drops before anyone suspected the unit.

The rule: \textbf{carry the finest unit the hardware gives you and divide only where a human reads
it}. This project's classifier now takes microseconds everywhere and converts once, at the boundary,
for display.

\subsection{The Cycle Counter: Measurement, Not Delay}

The third instrument is not for waiting at all:

\begin{minted}{rust}
let start = dwt::cycles();
// ... work ...
let elapsed_us = dwt::cycles_to_us(dwt::since(start), 19_000_000);
\end{minted}

Two registers and no ownership questions: enable \codeinline{TRCENA} in \codeinline{DEMCR}, then
\codeinline{CYCCNTENA} in \codeinline{DWT->CTRL}, then read \codeinline{CYCCNT}. It is how every
timing number in this book was obtained, and it is the only instrument here that cannot lie about
the clock because the CPU \emph{is} the clock.

One trap, measured: a counter that was never enabled reads as zero, and zero elapsed cycles looks
exactly like a loop that ran impossibly fast. Compare two consecutive reads before trusting either.
